\documentclass[11pt]{article, linenumbers}
\usepackage[margin=1.1in]{geometry}
\usepackage{parskip}
\pagestyle{empty}
\begin{document}
\noindent Ron Bibb\\
Independent Researcher\\
Lilburn, GA, USA\\
ronbibb@gmail.com \quad ORCID 0009-0004-1153-2464

\vspace{1em}
\noindent July 5, 2026

\vspace{1em}
\noindent Dear Editor,

Please consider the enclosed manuscript, \emph{``Residual Radial Structure in SPARC Dark
Matter Halos Reproduces on Independently Derived Rotation Curves,''} for publication in
\emph{The Astrophysical Journal Letters}.

Single smooth equilibrium halo profiles leave statistically significant, radially organized
residuals in ${\sim}35\%$ of SPARC rotation curves. Because such a detection could reflect
fitting-basis flexibility or a systematic of one velocity-field extraction rather than a
property of the galaxies, this Letter reports the most direct test of galaxy-intrinsic
origin: whether the model selection reproduces when rotation curves are re-derived from the
original radio data cubes by independent software, holding the baryonic model fixed. It
does, in four of five adjudicable systems --- decisively for NGC~2403
($\Delta\mathrm{BIC}=121$, on a curve derived from different observations) and NGC~3198
($\Delta\mathrm{BIC}=55$) --- with the crossover radius preserved; the one non-reproduction
and one power-limited case are reported explicitly. Because SPARC underpins a large body of
current work on dark matter halo structure, we believe this concise, single-result
verification is of immediate and broad interest, and is naturally suited to the Letters
format.

A companion paper presenting the full method-independence and robustness program is in
preparation and is cited in the manuscript as Bibb~2026b; the present Letter is fully
self-contained. The manuscript is not under consideration elsewhere, and no presubmission
inquiry was made. All code, pipeline configurations, and derived data products are openly
archived (Zenodo, DOI 10.5281/zenodo.21196421) and developed at
github.com/RonBibb/sparc-sphere. The warp-aware rotation curves for NGC~2841 and NGC~3521
were provided by J.~Sellwood and K.~Spekkens and are credited as private communication in
the text.

Thank you for your consideration.

\vspace{1.5em}
\noindent Sincerely,\\[0.5em]
Ron Bibb
\end{document}
